\chapter{فناوری‌ها و پیشینه}

این فصل فناوری‌های اصلی پروژه را معرفی می‌کند و توضیح می‌دهد هر کدام در این سامانه چه نقشی دارند. معیار انتخاب واژه‌ها هم ساده است: هر جا اصطلاح فارسی دقیق و جاافتاده وجود دارد از فارسی استفاده شده، و هر جا ترجمه باعث ابهام می‌شود، اصطلاح اصلی انگلیسی آمده است.

\section{\lr{Python}، \lr{Django} و \lr{Django REST Framework}}
\lr{Python} زبان اصلی بک‌اند است؛ زبانی که به خاطر خوانایی و اکوسیستم گسترده، برای توسعه وب و ساخت سرویس‌های داده‌محور زیاد استفاده می‌شود \cite{python}. روی این زبان، \lr{Django} چارچوب اصلی وب را فراهم می‌کند: تنظیمات، مسیرها، امنیت پایه، مدیریت کاربران، پنل مدیریت و ابزارهای آزمون \cite{django}. 

\lr{Django REST Framework} لایه \lr{API} را روی \lr{Django} کامل‌تر می‌کند. مفاهیمی مثل \lr{View}، \lr{Serializer}، اعتبارسنجی ورودی و پاسخ \lr{JSON} کمک می‌کنند مسیرهای عمومی و خصوصی سامانه با قرارداد منظم‌تری نوشته شوند \cite{drf}. در این پروژه، همین ترکیب برای ساخت کاتالوگ، حساب کاربری، اشتراک، مصرف و ابزارهای توسعه‌دهنده استفاده شده است.

\section{\lr{MongoDB} و مدل داده}
برای داده عملیاتی، پروژه به مدل سندی نزدیک است؛ مدلی که با ساختارهایی مانند کاربر، نشست، \lr{API}، پلن، مصرف و جریان \lr{Studio} راحت‌تر کار می‌کند. \lr{MongoDB} در داکیومنت رسمی خود همین رویکرد سندمحور را برای ذخیره داده‌های انعطاف‌پذیر توضیح می‌دهد \cite{mongodb}. در متن گزارش، وقتی از لایه داده حرف زده می‌شود، منظور همین بخش است که داده‌های محصول را برای مسیرهای بک‌اند آماده می‌کند.

این مدل برای بازارچه \lr{API} منطقی است، چون همه داده‌ها قالب یکسان ندارند. مستندات یک \lr{API}، پلن‌های قیمت‌گذاری، رخدادهای مصرف و جریان‌های \lr{Studio} هر کدام شکل خودشان را دارند. مدل سندی باعث می‌شود این تفاوت‌ها با انعطاف بیشتری نگهداری شوند.

\section{\lr{React}، \lr{TypeScript} و \lr{Vite}}
فرانت‌اند با \lr{React} ساخته شده است؛ کتابخانه‌ای که برای ساخت رابط‌های کاربری مبتنی بر کامپوننت استفاده می‌شود \cite{react}. \lr{TypeScript} هم به کد سمت کاربر نوع‌دهی اضافه می‌کند تا خطاهای ساده زودتر دیده شوند و قرارداد بین بخش‌های مختلف روشن‌تر بماند \cite{typescript}. 

\lr{Vite} ابزار توسعه و \lr{build} فرانت‌اند است. داکیومنت رسمی \lr{Vite} روی سرعت راه‌اندازی و تجربه توسعه سریع تاکید دارد \cite{vite}. این ویژگی برای پروژه‌ای با صفحه‌های مختلف مانند کاتالوگ، جزئیات \lr{API}، \lr{Dashboard}، \lr{Caller} و \lr{Studio} مفید است، چون توسعه رابط کاربری را روان‌تر می‌کند.

\section{\lr{Node.js} و ابزارهای فرانت‌اند}
اکوسیستم فرانت‌اند پروژه روی \lr{Node.js} و ابزارهای بسته‌بندی جاوااسکریپت می‌چرخد. \lr{Node.js} اجرای جاوااسکریپت خارج از مرورگر را فراهم می‌کند و پایه بسیاری از ابزارهای توسعه فرانت‌اند است \cite{nodejs}. کتابخانه‌هایی مثل \lr{TanStack Query}، \lr{Axios} و مجموعه‌های کامپوننتی، کار ارتباط با بک‌اند، مدیریت وضعیت درخواست‌ها و ساخت فرم‌ها و جدول‌ها را ساده‌تر می‌کنند.

در گزارش، نام این ابزارها بدون ترجمه آمده است، چون ترجمه فارسی آن‌ها معمولاً دقیق نیست و ممکن است خواننده را از اصطلاح رایج دور کند.

\section{\lr{Docker Compose} و اجرای محلی}
\lr{Docker Compose} برای تعریف و اجرای چند سرویس کنار هم استفاده می‌شود \cite{docker}. در این پروژه، همین فضا باعث می‌شود بک‌اند و فرانت‌اند در کنار هم بالا بیایند و مسیر ارتباطی مشخصی داشته باشند. این روش برای ارائه، آزمون دستی و توسعه محلی مناسب است، هرچند جای پیکربندی تولیدی را نمی‌گیرد.

\section{\lr{OpenAPI} و \lr{RapidAPI}}
\lr{OpenAPI} برای توصیف قرارداد \lr{API} به‌کار می‌رود؛ یعنی مسیرها، ورودی‌ها، خروجی‌ها و پاسخ‌ها را به شکلی می‌نویسد که هم انسان و هم ابزار بتوانند آن را بخوانند \cite{openapi}. \lr{RapidAPI} نیز نمونه شناخته‌شده‌ای از فضای بازارچه \lr{API} است و برای فهم تجربه کاربر در کشف و مصرف \lr{API} مرجع خوبی به شمار می‌آید \cite{rapidapi}.

\section{جمع‌بندی فصل}
پشته فنی پروژه با ماهیت آن هماهنگ است: \lr{Django REST Framework} برای ساخت \lr{API}، \lr{React} برای تجربه کاربری، \lr{MongoDB} برای داده‌های سندی، \lr{Docker Compose} برای اجرای محلی، و \lr{OpenAPI} برای خوانایی قرارداد سرویس. این ترکیب، فضای گزارش را هم به سمت یک محصول واقعی اشتراک‌گذاری \lr{API} می‌برد.
